\section{Conclusion}
\label{sec:conclusion}

We audited a released fNIRS foundation model on a clean held-out dataset. The pretrained
encoder reached $0.469 \pm 0.095$ three-class accuracy against chance $0.333$, but a randomly
initialised encoder of the same architecture reached $0.448 \pm 0.010$ over ten seeds---a
$2.0$~pp difference that is not statistically significant. The architecture itself did
significantly outperform hand-crafted features, and that advantage disappeared once
short-separation channels were regressed out, at a cost of $4.5$~pp. Naive concatenation of
representations significantly degraded performance.

The methodological conclusion is modest and, we think, portable: fNIRS foundation-model claims
should be reported against a random-encoder baseline, with contamination accounting and with
preprocessing sensitivity. The empirical conclusion is deliberately narrow: on this dataset and
task, the pretraining advantage of a fNIRS foundation model is not established.

Three extensions follow naturally. First, replicate across datasets and paradigms, including a
visual or affective domain the model's authors never tested. Second, fine-tune rather than
probe, and evaluate under the same controls. Third, replace naive concatenation with learned
fusion, and test whether it beats the better single representation. We release the harness, the
frozen protocol and a detectability calculator so that all three can be attempted without
rebuilding the infrastructure.

\section*{Data and code availability}
\label{sec:availability}
All analysis code, the frozen evaluation protocol, per-fold results and figure-generation scripts
are available at \url{https://github.com/curiousbrutus/fnirs-foundation-model-audit} and archived
on Zenodo (DOI to be added on acceptance). The auditory dataset analysed in
\S\ref{sec:results} is publicly available from its original authors \citep{luke2021passive}.

The trigger-free pilot recordings used only to illustrate the detectability calculator
(\S\ref{sec:results_tool}) are \emph{not} released. They are informal, self-collected exploratory
recordings made by the author, obtained without a formal ethics protocol or research consent, and
they are reported exclusively as aggregate statistics under neutral session codes.

\section*{Ethics statement}
This study re-analyses previously collected, publicly available, de-identified data; no new human
participants were involved and no new ethical approval was required.

\section*{Author contributions}
E.G. conceived the study, performed the analysis and wrote the manuscript.
\texttt{[TODO: add co-authors and CRediT roles before submission.]}

\section*{Use of AI tools}
Generative AI tools were used in the preparation of this work: a coding assistant was used to
implement, debug and test the analysis pipeline, and a language model was used to draft and edit
prose. The research question, the pre-registered evaluation protocol, the statistical choices and
all interpretations are the author's. Every number reported here was recomputed by the author
from the released code and result files, and the author takes full responsibility for the content
of the manuscript.

\section*{Acknowledgment}
The author thanks Zenghui Wang (Southeast University, Nanjing, China) for generously sharing his
pretrained fNIRS models and code, and for his body of work on fNIRS classification architectures
\citep{wang2022vision,wang2022transformer,wang2023rethinking,wang2026unified,wang2026dkssm};
the authors of \citet{luke2021passive} for releasing their data; and the MtReNIRS authors for
releasing model weights under a permissive licence, which made an independent audit possible.
